% thesis_diabetes_gradcam.tex -- Report for Explainable AI: Grad-CAM for Diabetes CNN
% Replace content in your existing thesis.tex or merge with dalthesis template

\documentclass[12pt]{dalthesis}

% Page layout and geometry
\usepackage{geometry}
\geometry{a4paper, margin=1in}

% Table and array packages
\usepackage{longtable}
\usepackage{array}
\usepackage{nomencl}
\usepackage{booktabs}
\usepackage{multirow}
\usepackage{tabularx}
\usepackage{adjustbox}
\usepackage{nomencl}
\makenomenclature

% Graphics and figures
\usepackage{graphicx}
\usepackage{caption}
\usepackage{float}
\usepackage{pdflscape}

% Mathematical packages
\usepackage{amsmath}
\usepackage{amssymb}

% Algorithms
\usepackage{algorithm}
\usepackage{algorithmic}

% Hyperlinks
\usepackage{hyperref}

% Multi-column and floating figures
\usepackage{multicol}
\usepackage{placeins}

% Code listings
\usepackage{listings}
\usepackage{xcolor}
\usepackage{colortbl}
\usepackage{pifont}

\definecolor{codegreen}{rgb}{0,0.6,0}
\definecolor{codegray}{rgb}{0.5,0.5,0.5}
\lstset{
    basicstyle=\ttfamily\small,
    breaklines=true,
    frame=single,
    numbers=left,
    numberstyle=\tiny,
    keywordstyle=\color{blue},
    commentstyle=\color{codegreen},
}

\newcommand{\xmark}{\ding{55}}

\begin{document}

\title{Explainable AI for Diabetes Risk Prediction Using Grad-CAM and IGTD}
\author{Name of the student}

\macs  
\degreeinitial{M.A.C.Sc.}
\faculty{Computer Science}
\dept{Faculty of Computer Science}

\defencemonth{Dec}
\defenceyear{2024}

\makenomenclature
\renewcommand{\nomname}{List of Abbreviations and Symbols Used}

\nomenclature[A]{\textbf{CDC}}{Centers for Disease Control and Prevention}
\nomenclature[A]{\textbf{CNN}}{Convolutional Neural Network}
\nomenclature[A]{\textbf{ENN}}{Edited Nearest Neighbors}
\nomenclature[A]{\textbf{Grad-CAM}}{Gradient-weighted Class Activation Mapping}
\nomenclature[A]{\textbf{IGTD}}{Image Generator for Tabular Data}
\nomenclature[A]{\textbf{XAI}}{Explainable Artificial Intelligence}

\frontmatter

\begin{abstract}
This report presents an Explainable AI (XAI) approach for diabetes risk prediction using a Convolutional Neural Network (CNN) combined with the Image Generator for Tabular Data (IGTD) algorithm. We transform tabular health data from the CDC diabetes dataset into 2D grayscale images, train a CNN based on Shenghao Wang et al.'s cnn\_igtd\_f15\_enn model, and apply Grad-CAM to visualize which features the model focuses on when making predictions. The pipeline achieves 84\% accuracy on the test set and produces interpretable heatmaps with feature names and color scales, supporting trust and transparency in the model's decisions. The implementation is available on GitHub: \url{https://github.com/mzaman2016/VitaScreen}.
\end{abstract}

\addcontentsline{toc}{chapter}{List of Abbreviations and Symbols Used}
\printnomenclature

\begin{acknowledgements}
% Add your acknowledgements here.
\end{acknowledgements}

\mainmatter

\chapter{Introduction}
\section{Overview}

Diabetes is a major global health concern, and early risk assessment can support preventive care. Machine learning models can predict diabetes risk from health indicators, but black-box models lack interpretability. This work applies Explainable AI (XAI) using Grad-CAM to a CNN that predicts diabetes from tabular data transformed into images via IGTD.

\section{Problem Statement}

Deep learning models for healthcare often lack transparency. Clinicians and patients need to understand \textit{why} a model predicts a certain outcome. We address this by integrating Grad-CAM into a diabetes prediction pipeline to highlight which input features drive each prediction.

\section{Motivation}

Trust in AI systems requires explainability. Grad-CAM provides visual explanations by showing which regions of the input (IGTD images) the CNN focuses on. This helps validate model behavior and identify potential biases.

\section{Major Contributions}

\begin{itemize}
\item Implemented a full pipeline: CDC dataset $\rightarrow$ IGTD $\rightarrow$ CNN $\rightarrow$ Grad-CAM.
\item Added feature names and color scales to Grad-CAM visualizations for interpretability.
\item Achieved 84\% accuracy on the test set with the cnn\_igtd\_f15\_enn architecture.
\item Integrated with the VitaScreen GitHub repository.
\end{itemize}

\section{Organization of the report}

Chapter~\ref{ch:background} covers background on IGTD, CNNs, and Grad-CAM. Chapter~\ref{ch:lit} reviews related work. Chapter~\ref{ch:method} describes the methodology. Chapter~\ref{ch:impl} presents implementation details. Chapter~\ref{ch:results} discusses results. Chapter~\ref{ch:conclusion} concludes with limitations and future work.

\chapter{Background Knowledge}\label{ch:background}

\section{IGTD (Image Generator for Tabular Data)}

IGTD converts tabular data into 2D images by optimizing the arrangement of features in a grid so that similar features are placed close together. We use a 5$\times$3 grid for 15 features.

\section{CNN Architecture}

We use Shenghao Wang et al.'s cnn\_igtd\_f15\_enn: 4 convolutional blocks (64 filters, 3$\times$3), BatchNorm, ReLU, Global Average Pooling, Dropout (50\%), and a fully connected output layer.

\section{Grad-CAM}

Grad-CAM computes a heatmap by combining the gradients of the target class with respect to the last convolutional layer's activations. Red indicates high importance; blue indicates low importance.

\chapter{Literature Review}\label{ch:lit}

Zhu et al.~\cite{zhu2021igtd} introduced IGTD for converting tabular data into images. Wang et al. applied IGTD to diabetes prediction with CNNs. Selvaraju et al.~\cite{selvaraju2017gradcam} proposed Grad-CAM for CNN interpretability. The VitaScreen project~\cite{vitascreen} combines IGTD and NCTD with deep learning for diabetes risk prediction.

\chapter{Proposed Methodology}\label{ch:method}

\section{Pipeline Overview}

\begin{enumerate}
\item Load CDC diabetes dataset (or cdcNormalDiabeticFE1\_20RFFSQ.csv).
\item Select 15 features: HighBP, HighChol, CholCheck, BMI, Smoker, Stroke, HeartDiseaseorAttack, PhysActivity, Fruits, Veggies, HvyAlcoholConsump, AnyHealthcare, NoDocbcCost, GenHlth, MentHlth.
\item Split 80:20 train/test; apply ENN (k=3) on training set.
\item Run IGTD to create 5$\times$3 grayscale images.
\item Train CNN for 30 epochs.
\item Apply Grad-CAM on test samples; save heatmaps with feature names and scales.
\end{enumerate}

\section{Grad-CAM with Feature Labels}

We overlay feature names on each cell of the IGTD image and Grad-CAM heatmap. Two color scales are used: (1) IGTD: White=Low, Black=High feature value; (2) Grad-CAM: Red=High, Blue=Low importance.

\chapter{Implementation}\label{ch:impl}

\section{Before Grad-CAM: Base Pipeline}

The base pipeline (\texttt{diabetes\_prediction\_pipeline.py}) loads data, runs IGTD, trains the CNN, and evaluates. It does not include explainability.

\section{After Grad-CAM: XAI Pipeline}

The XAI pipeline (\texttt{gradcam\_explainability.py}) extends the base pipeline with:
\begin{itemize}
\item Grad-CAM class that hooks the last convolutional layer.
\item Feature-to-position mapping from IGTD index.
\item Visualization with feature names and color scales on both IGTD and heatmap panels.
\end{itemize}

\section{Code Structure}

\begin{itemize}
\item \texttt{diabetes\_prediction\_pipeline.py}: Base CNN + IGTD pipeline.
\item \texttt{gradcam\_explainability.py}: Full pipeline with Grad-CAM.
\item \texttt{IGTD\_Functions.py}: IGTD algorithm (from Zhu et al.).
\end{itemize}

\section{GitHub Repository}

The project is integrated with VitaScreen: \url{https://github.com/mzaman2016/VitaScreen}. Clone the repository and run \texttt{python gradcam\_explainability.py} to reproduce results.

\chapter{Result Discussion}\label{ch:results}

\section{Performance Metrics}

\begin{table}[H]
\centering
\caption{Test Set Performance}
\begin{tabular}{@{}lll@{}}
\toprule
\textbf{Metric} & \textbf{Value} & \textbf{Description} \\
\midrule
Accuracy  & 0.84 & 84\% of all predictions are correct \\
Precision & 0.38 & When the model predicts diabetic, it is right 38\% of the time \\
Recall    & 0.50 & The model detects 50\% of all diabetic cases \\
F1-Score  & 0.43 & Balance between Precision and Recall \\
\bottomrule
\end{tabular}
\end{table}

\section{Classification Report}

\begin{table}[H]
\centering
\caption{Per-Class Performance}
\begin{tabular}{@{}lrrrrr@{}}
\toprule
\textbf{Class} & \textbf{Precision} & \textbf{Recall} & \textbf{F1-Score} & \textbf{Support} \\
\midrule
Non-Diabetic & 0.93 & 0.89 & 0.91 & 88 \\
Diabetic    & 0.38 & 0.50 & 0.43 & 12 \\
\midrule
Accuracy    & \multicolumn{4}{r}{0.84} & 100 \\
\bottomrule
\end{tabular}
\end{table}

\section{Grad-CAM Visualizations}

The Grad-CAM outputs (saved in \texttt{FDI/}) show:
\begin{itemize}
\item \textbf{IGTD Image}: Grayscale 5$\times$3 grid with feature names; scale: White=Low, Black=High.
\item \textbf{Grad-CAM Heatmap}: Color-coded importance; scale: Red=High, Blue=Low.
\end{itemize}

\textbf{Figures to attach:}
\begin{itemize}
\item Figure~\ref{fig:gradcam1}: Sample 1 -- Non-Diabetic, correctly predicted.
\item Figure~\ref{fig:gradcam2}: Sample 2 -- Diabetic, incorrectly predicted as Non-Diabetic.
\item Figure~\ref{fig:gradcam3}: Sample 3 -- Non-Diabetic, correctly predicted.
\item Figure~\ref{fig:gradcam4}: Sample 4 -- Non-Diabetic, correctly predicted.
\item Figure~\ref{fig:gradcam5}: Sample 5 -- Non-Diabetic, correctly predicted.
\item Figure~\ref{fig:gradcam6}: Sample 6 -- Non-Diabetic, correctly predicted.
\item Figure~\ref{fig:gradcam7}: Sample 7 -- Non-Diabetic, correctly predicted.
\item Figure~\ref{fig:gradcam8}: Sample 8 -- Non-Diabetic, correctly predicted.
\end{itemize}

% Create figures/ folder and copy gradcam_sample_1.png through gradcam_sample_8.png from FDI/
\begin{figure}[H]
\centering
\includegraphics[width=0.9\textwidth]{figures/gradcam_sample_1.png}
\caption{Grad-CAM Sample 1: True=0 (Non-Diabetic), Pred=0. IGTD image (left) and Grad-CAM heatmap (right) with feature names and scales.}
\label{fig:gradcam1}
\end{figure}

\begin{figure}[H]
\centering
\includegraphics[width=0.9\textwidth]{figures/gradcam_sample_2.png}
\caption{Grad-CAM Sample 2: True=1 (Diabetic), Pred=0 -- Misclassified case.}
\label{fig:gradcam2}
\end{figure}

\begin{figure}[H]
\centering
\includegraphics[width=0.9\textwidth]{figures/gradcam_sample_3.png}
\caption{Grad-CAM Sample 3: True=0 (Non-Diabetic), Pred=0.}
\label{fig:gradcam3}
\end{figure}

\begin{figure}[H]
\centering
\includegraphics[width=0.9\textwidth]{figures/gradcam_sample_4.png}
\caption{Grad-CAM Sample 4: True=0 (Non-Diabetic), Pred=0.}
\label{fig:gradcam4}
\end{figure}

\begin{figure}[H]
\centering
\includegraphics[width=0.9\textwidth]{figures/gradcam_sample_5.png}
\caption{Grad-CAM Sample 5: True=0 (Non-Diabetic), Pred=0.}
\label{fig:gradcam5}
\end{figure}

\begin{figure}[H]
\centering
\includegraphics[width=0.9\textwidth]{figures/gradcam_sample_6.png}
\caption{Grad-CAM Sample 6: True=0 (Non-Diabetic), Pred=0.}
\label{fig:gradcam6}
\end{figure}

\begin{figure}[H]
\centering
\includegraphics[width=0.9\textwidth]{figures/gradcam_sample_7.png}
\caption{Grad-CAM Sample 7: True=0 (Non-Diabetic), Pred=0.}
\label{fig:gradcam7}
\end{figure}

\begin{figure}[H]
\centering
\includegraphics[width=0.9\textwidth]{figures/gradcam_sample_8.png}
\caption{Grad-CAM Sample 8: True=0 (Non-Diabetic), Pred=0.}
\label{fig:gradcam8}
\end{figure}

\section{Key Findings}

\begin{itemize}
\item The model performs better on non-diabetic cases (88 samples) than diabetic cases (12 samples) due to class imbalance.
\item Grad-CAM highlights features such as Veggies, AnyHealthcare, CholCheck, and MentHlth as influential for non-diabetic predictions.
\item Feature names and scales improve interpretability for clinicians and researchers.
\end{itemize}

\chapter{Conclusion}\label{ch:conclusion}

We implemented an Explainable AI pipeline for diabetes risk prediction using Grad-CAM with IGTD-transformed tabular data. The model achieves 84\% accuracy, and Grad-CAM heatmaps with feature names and scales provide transparent explanations of model decisions. The work is integrated with the VitaScreen GitHub repository.

\chapter{Limitations and Future Work}
\section{Limitations}

\begin{itemize}
\item Class imbalance (88 non-diabetic vs.\ 12 diabetic in test set) affects diabetic-class performance.
\item Grad-CAM explains spatial regions; direct feature importance may require additional methods (e.g., SHAP).
\item Dataset used 500 samples; full CDC dataset may yield different results.
\end{itemize}

\section{Future Work}

\begin{itemize}
\item Apply SHAP or LIME for tabular-level feature importance.
\item Use the full CDC dataset or cdcNormalDiabeticFE1\_20RFFSQ.csv when available.
\item Integrate Grad-CAM into a web-based VitaScreen prototype.
\item Explore NCTD as an alternative to IGTD for tabular-to-image conversion.
\end{itemize}

\bibliographystyle{plain}
\bibliography{report}

\end{document}
